\subsection{Rangos de edad}

Tenemos reglas que nos indican que los rangos de edad más comunes entre las víctimas desaparecidas.

\begin{table}[H]
\centering
\small
\begin{tabular}{cccccc}
\hline
ID & Antecedente & Consecuente & Soporte & Confianza & Lift \\
\hline
24 & \texttt{GRUPO\_EDAD=$50+$} & \texttt{ESTATUS\_VICTIMA=DESAPARECIDA} & 0.0840 & 0.9314 & 1.5611 \\
25 & \texttt{GRUPO\_EDAD=31-50} & \texttt{ESTATUS\_VICTIMA=DESAPARECIDA} & 0.1881 & 0.3154 & 1.5723 \\
26 & \texttt{ESTATUS\_VICTIMA=DESAPARECIDA} & \texttt{GRUPO\_EDAD=31-50} & 0.1882 & 0.3154 & 1.5723 \\
27 & \texttt{GRUPO\_EDAD=18-30} & \texttt{ESTATUS\_VICTIMA=DESAPARECIDA} & 0.0874 & 0.9348 & 0.9348 \\
\hline
\end{tabular}
\caption{Regla directa para rango de edad condicionado a estatus desaparecida. Valores redondeados.}
\label{tab:edad_dado_desaparecida}
\end{table}

Estas reglas sugieren que el rango de edad promedio de las víctimas desaparecidas es entre 31 y 50 años (con valores de
confianza de 0.3154). Esto puede indicar que las personas en este rango de edad son más vulnerables a la desaparición, o
que hay un sesgo en el registro hacia estos grupos demográficos. Sin embargo, también es importante considerar que el rango
de edad de 18-30 años tiene una confianza más alta (0.9348) para el estatus de desaparecida, lo que podría indicar que aunque
hay menos casos en este rango, es más probable que sean reportados como desaparecidos.

\subsection{Perfil Demográfico}

Las siguientes reglas sugieren un perfil demográfico: conforme aumenta la edad, se observa una mayor sobrerrepresentación de hombres en el registro:

\begin{table}[H]
\centering
\small
\begin{tabular}{cccccc}
\hline
ID & Antecedente & Consecuente & Soporte & Confianza & Lift \\
\hline
14 & \texttt{GRUPO\_EDAD=50+} & \texttt{SEXO=HOMBRE} & 0.0790 & 0.8756 & 1.813 \\
15 & \texttt{GRUPO\_EDAD=31-50} & \texttt{SEXO=HOMBRE} & 0.1671 & 0.8330 & 1.725 \\
17 & \texttt{GRUPO\_EDAD=18-30} & \texttt{SEXO=HOMBRE} & 0.0634 & 0.6781 & 1.404 \\
53 & \texttt{SEXO=HOMBRE, ESTATUS\_VICTIMA=DESAPARECIDA} & \texttt{GRUPO\_EDAD=31-50} & 0.1567 & 0.3449 & 1.719 \\
\hline
\end{tabular}
\caption{Reglas de asociación para perfil demográfico (edad y sexo). Valores redondeados.}
\label{tab:perfil_demografico}
\end{table}

Aquí, los lifts crecientes ($<1$) por grupo sugieren un sesgo de reporte/captura por edad y sexo.

\subsection{Confidencialidad en Bloque}

Estas reglas muestran un patrón sistémico: cuando un campo es confidencial, los demás también lo son, con una confianza
de 1.0, en una proporción muy grande del dataset.

\begin{table}[htbp]
\centering
\small
\begin{tabular}{cccccc}
\hline
ID & Antecedente & Consecuente & Support & Conf. & Lift \\
\hline
7  & \texttt{ESTATUS\_VICTIMA=CONFIDENCIAL} & \texttt{SEXO=CONFIDENCIAL} & 0.3671 & 1.0000 & 2.724 \\
10 & \texttt{SEXO=CONFIDENCIAL} & \texttt{MUNICIPIO=CONFIDENCIAL} & 0.3671 & 1.0000 & 2.724 \\
20 & \texttt{ESTATUS\_VICTIMA=CONFIDENCIAL} & \texttt{MUNICIPIO=CONFIDENCIAL} & 0.3671 & 1.0000 & 2.724 \\
\hline
\end{tabular}
\caption{Sesgo de registro: confidencialidad en bloque (censura/datos faltantes estructurados). Valores redondeados.}
\label{tab:confidencialidad_bloque}
\end{table}

Esto nos lleva a esta conclusión: existe un mecanismo de censura sistémico de la información. Se realiza el registro
de la persona desaparecida, pero no se deciden hacer públicos ninguno de los datos personales. Esto puede deberse a
políticas de privacidad (sin explicar cuales son), temor a represalias, o falta de información capturada en el momento
del registro.

\subsection{Sesgo por entidad — caso Jalisco}

Estas reglas son un triste ejemplo de lo anterior. Si el registro de desaparición se hace en Jalisco, es muy probable
que el resto de los campos terminen siendo confidenciales, lo que impide cualquier análisis de perfil o tendencia en esta entidad.

\begin{table}[H]
\centering
\small
\begin{tabular}{cccccc}
\hline
ID & Antecedente & Consecuente & Soporte & Confianza & Lift \\
\hline
9  & \texttt{ENTIDAD=JALISCO} & \texttt{SEXO=CONFIDENCIAL} & 0.0747 & 0.6752 & 1.839 \\
19 & \texttt{ENTIDAD=JALISCO} & \texttt{ESTATUS\_VICTIMA=CONFIDENCIAL} & 0.0747 & 0.6752 & 1.839 \\
30 & \texttt{ENTIDAD=JALISCO} & \texttt{MUNICIPIO=CONFIDENCIAL} & 0.0747 & 0.6752 & 1.839 \\
36 & \texttt{ESTATUS\_VICTIMA=CONFIDENCIAL, ENTIDAD=JALISCO} & \texttt{SEXO=CONFIDENCIAL} & 0.0747 & 1.0000 & 2.724 \\
45 & \texttt{ENTIDAD=JALISCO', 'SEXO=CONFIDENCIAL} & \texttt{MUNICIPIO=CONFIDENCIAL} & 0.0747 & 1.0000 & 2.724 \\
46 & \texttt{ENTIDAD=JALISCO', 'MUNICIPIO=CONFIDENCIAL} & \texttt{SEXO=CONFIDENCIAL} & 0.0747 & 1.0000 & 2.724 \\
\hline
\end{tabular}
\caption{Sesgo por entidad: patrones de confidencialidad asociados a Jalisco. Valores redondeados.}
\label{tab:jalisco_confidencialidad}
\end{table}

\subsection{Sesgo Temporal}

En base a estas dos reglas podemos inferir que los registros de 2024 en adelante
tienen un rezago de actualización, ya que no existe existe un cambio en el estado
de la víctima. Esto se puede deber a que el proceso de investigación y actualización
de casos es lento, o a que los casos más recientes aún no han sido actualizados en el
sistema.

\begin{table}[H]
\centering
\small
\begin{tabular}{cccccc}
\hline
ID & Antecedente & Consecuente & Soporte & Confianza & Lift \\
\hline
28 & \texttt{AÑO\_DESAPARICION=2025.0} & \texttt{ESTATUS\_VICTIMA=DESAPARECIDA} & 0.0512 & 0.9985 & 1.674 \\
29 & \texttt{AÑO\_DESAPARICION=2024.0} & \texttt{ESTATUS\_VICTIMA=DESAPARECIDA} & 0.0598 & 0.9960 & 1.669 \\
\end{tabular}
\caption{Sesgo temporal: asociación entre año reciente y estatus ``DESAPARECIDA'' (posible rezago de actualización). Valores redondeados.}
\label{tab:sesgo_temporal}
\end{table}

Los archivos con los cuales se hicieron estas interpretaciones se pueden encontrar
aqui:
\begin{itemize}
    \item \href{https://github.com/Yoba7-7/AyMdD_2026-2/tree/main/Tarea%203/datasets#:~:text=datos_manual_final.csv}{Reglas Finales con nuestra
    implementación de Apriori.}
    \item \href{https://github.com/Yoba7-7/AyMdD_2026-2/blob/main/Tarea%203/datasets/datos_mlxtend_final.csv}{Reglas Finales con la implementación de mlxtend.}
\end{itemize}